\documentclass[11pt]{article}
\usepackage{graphicx}
\usepackage{booktabs}
\usepackage[margin=22mm]{geometry}
\title{CV-MsAtViT 3地点レポート}
\date{}
\begin{document}
\maketitle
\begin{table}[h]
  \centering
  \begin{tabular}{lrrr}
    \toprule
    データセット & OA(%) & AA(%) & カッパ×100 \ 
    \midrule
    Flevoland & 97.45 & 96.76 & 97.22 \
    San Francisco & 97.86 & 94.53 & 96.65 \
    Oberpfaffenhofen & 59.90 & 45.64 & 18.85 \
    \bottomrule
  \end{tabular}
\end{table}
\section{Flevoland}
\begin{figure}[h]
  \centering
  \includegraphics[width=0.95\linewidth]{results/paper/figs/triptych_Full_FL_T.png}
  \caption{Flevoland: PauliRGB/GT/予測}
\end{figure}

\begin{table}[h]
  \centering
  \begin{tabular}{lrrr}
    \toprule
    指標 & OA(%) & AA(%) & カッパ×100 \
    \midrule
    値 & 97.45 & 96.76 & 97.22 \
    \bottomrule
  \end{tabular}
\end{table}
\section{San Francisco}
\begin{figure}[h]
  \centering
  \includegraphics[width=0.95\linewidth]{results/paper/figs/triptych_Full_SF.png}
  \caption{San Francisco: PauliRGB/GT/予測}
\end{figure}

\begin{table}[h]
  \centering
  \begin{tabular}{lrrr}
    \toprule
    指標 & OA(%) & AA(%) & カッパ×100 \
    \midrule
    値 & 97.86 & 94.53 & 96.65 \
    \bottomrule
  \end{tabular}
\end{table}
\section{Oberpfaffenhofen}
\begin{figure}[h]
  \centering
  \includegraphics[width=0.95\linewidth]{results/paper/figs/triptych_Full_ober.png}
  \caption{Oberpfaffenhofen: PauliRGB/GT/予測}
\end{figure}

\begin{table}[h]
  \centering
  \begin{tabular}{lrrr}
    \toprule
    指標 & OA(%) & AA(%) & カッパ×100 \
    \midrule
    値 & 59.90 & 45.64 & 18.85 \
    \bottomrule
  \end{tabular}
\end{table}
\end{document}